\VerdierPhysicalPageAnchor{0072}%
\VerdierSemanticAnchor{ch1-2-3-2-proof}
An immediate calculation shows that the family
$\bigl((p_{i,\varphi}(Y))^\tau,\ \tau\in\mathbb Z[m]\bigr)$ commutes
with the differentials of $\int_\varphi^\Sigma T_iY$ and of
$T_{\varphi(i)}\int_\varphi^\Sigma Y$, and consequently defines a
morphism:

\VerdierSemanticAnchor{ch1-2-3-2-eq2}
\begin{equation}
  p_{i,\varphi}(Y)\colon
  \int_\varphi^\Sigma T_iY
  \longrightarrow
  T_{\varphi(i)}\int_\varphi^\Sigma Y.
  \tag{2.3.2.2}
\end{equation}

\VerdierSemanticAnchor{ch1-2-3-2-functorial-isomorphism}
This morphism is a functorial isomorphism in $Y$, whence an isomorphism
of functors:

\VerdierSemanticAnchor{ch1-2-3-2-eq3}
\begin{equation}
  p_{i,\varphi}\colon
  \int_\varphi^\Sigma T_i
  \xrightarrow{\ \sim\ }
  T_{\varphi(i)}\int_\varphi^\Sigma.
  \tag{2.3.2.3}
\end{equation}

\VerdierSemanticAnchor{ch1-2-3-3}
\noindent\textbf{Proposition 2.3.3.}\quad
\emph{1) For every pair $i,j\in[n]$, $i\ne j$, the diagram:}

\VerdierSemanticAnchor{ch1-2-3-3-eq1}
\begin{equation}
  \resizebox{\textwidth}{!}{%
    $\begin{tikzcd}[ampersand replacement=\&,column sep=4.2em,row sep=4.8em]
      \displaystyle\int_\varphi^\Sigma T_iT_j
        \arrow[r,"{p_{i,\varphi}\star T_j}"]
        \arrow[d,"{p_{j,\varphi}\star T_i}"']
        \& \displaystyle T_{\varphi(i)}\int_\varphi^\Sigma T_j
          \arrow[d,"{T_{\varphi(i)}\star p_{j,\varphi}}"] \\
      \displaystyle T_{\varphi(j)}\int_\varphi^\Sigma T_i
        \arrow[r,"{T_{\varphi(j)}\star p_{i,\varphi}}"]
        \& \displaystyle T_{\varphi(i)}T_{\varphi(j)}\int_\varphi^\Sigma
    \end{tikzcd}$%
  }
  \tag{2.3.3.1}
\end{equation}

\VerdierSemanticAnchor{ch1-2-3-3-commutation-dichotomy}
\emph{is commutative when $\varphi(i)\ne\varphi(j)$ and
anticommutative when $\varphi(i)=\varphi(j)$.}

\VerdierSemanticAnchor{ch1-2-3-3-composition}
\emph{2) Let $\psi\colon[m]\longrightarrow[r]$ be a map. The diagram:}

\VerdierSemanticAnchor{ch1-2-3-3-composition-diagram}
\begin{equation*}
  \resizebox{\textwidth}{!}{%
    $\begin{tikzcd}[ampersand replacement=\&,column sep=2.8em,row sep=5.0em]
      \displaystyle\int_\psi^\Sigma\int_\varphi^\Sigma T_i
        \arrow[rr,"{\int_\psi^\Sigma\star p_{i,\varphi}}"]
        \arrow[d,equal]
        \& {}
        \& \displaystyle\int_\psi^\Sigma T_{\varphi(i)}\int_\varphi^\Sigma
          \arrow[d,"{p_{\varphi(i),\psi}\star\int_\varphi^\Sigma}"] \\
      \displaystyle\int_{\psi\varphi}^\Sigma T_i
        \arrow[r,"{p_{i,\psi\varphi}}"']
        \& \displaystyle T_{\psi\varphi(i)}\int_{\psi\varphi}^\Sigma
          \arrow[r,equal]
        \& \displaystyle T_{\psi\varphi(i)}\int_\psi^\Sigma\int_\varphi^\Sigma
    \end{tikzcd}$%
  }
\end{equation*}

\VerdierSemanticAnchor{ch1-2-3-3-canonical-isomorphism-caveat}
\emph{is commutative.} (The reader who does not wish to cheat will replace
the sign ``$=$'' by canonical isomorphisms.)

\VerdierSemanticAnchor{ch1-2-3-3-proof}
The proof is immediate and left to the reader. The first assertion
follows from (2.2.1.2).
